% Pursuing Stacks (cotes 134-2 to 134-8): archive cote 134-2, site page 29 (PDF page 30).
% letter to Quillen. Draft compared: ps1.tex:723-786.
% AI-transcribed by Claude Opus 5.5 (claude-opus-5-5, adaptive thinking, effort high), Message Batches API, batch ps1 (msgbatch_01Rv5hCaNTtUt4JHx1yqoHPL), 2026-11-12; not yet checked by eye.
\scanpage{134-2/29}
\mspage{\unsure{\#5|8'}}

The expression ``deduce\unsure{s|d} from'' in B)  means that we are adding arrows
$h: K \to F_i$ (each with preassigned source and target in $\underline{C}_n$),  with as
``new axioms'' on \inserted{the bunch of} these uniquely the two relations \deleted{\illegible{1}}\inserted{(1),(2)} of \over{\illegible{1}}{A}), the category
$\underline{C}_{n+1}$  being deduced from $\underline{C}_n$ in an obvious way,  as the solution of a uni-%
versal problem with\inserted{in} the category of all categories where binary amalgama-%
ted products exist,  and ``maps'' between these being functors which commute
\margin{this has to be corrected -- amalgamated sums exist in \unsure{$C_\infty$} only, and these sho\unsure{uld} be respected}
to those fibered products\inserted{(14)}. \deleted{In}\inserted{In} practical terms,  the arrows of $\underline{C}_{n+1}$ are
those deduced from the arrows in $\underline{C}_n$ and the ``new'' arrows $h$,  by combining
formal operations of composing arrows by $v\circ u$ ,  and taking \inserted{(binary)} amalgamated
products of arrows. \deleted{\illegible{1}}

NB  Of course the condition (1) in A) is necessary for the existence
\margin{$\|$}of an $h$ satisfying (2).  That it is sufficient too \deleted{can}\inserted{may} be viewed as an
extremely strong,\xout{xxxxxxx} ``universal'' version of coherence conditions,
concerning the various structures introduced on a $\infty$-groupoid. Intuitive-%
ly,  it means that whenever we have \deleted{\over{\illegible{1}}{two} ways of}\inserted{two ways of} associat\unsure{i}ng to a \inserted{finite} family \deleted{\illegible{1}}
$(u_i)_{i\in I}$  of objects of an $\infty$-groupoid,  $u_i \in F_{n(i)}$ ,  \xout{xxxxxxxxxx} subjected
to a \hand{``}standar\over{t}{d}\hand{''} set of relations on the $u_i$'s, an element of some $F_n$ ,  in
terms of the $\infty$-groupoid structure only,  then we have automatically
\inserted[typed]{a ``homotopy''}\xout{xxxxxxxxxxx} between these built-in in the very structure of the $\infty$-%
groupoid,  provided it makes at all a sense to ask for one (namely provi-%
ded condition (1) holds if $n \geq 1$).   I have the feeling moreover that
conditions (A) and (B) (plus the relation $\underline{C}_\infty = \varinjlim \underline{C}_n$ ) is all what will
be ever needed,  when using the definition of the structure species,
- plus of course the description of $\underline{C}_o$ , and the \inserted{implicit} \unsure{fact} that the categories
\margin{inaccurate -- see above}
$\underline{C}_n$ are stable under binary fiber products and the inclusion functors com-%
mute to these\inserted{(13)}. Of course,  the category which really interests us is
$\underline{C}_\infty$,  the description of the intermediate $\underline{C}_n$'s\hand{$|$} is merely tec\over{\illegible{1}}{h}nical \deleted{txt} -
the main point is that there should exist an increasing sequence $(C_n)$ of
subcategories of $\underline{C}_\infty$ ,  having same objects\over{,}{(}and ``same'' fiberproducts\over{,}{)}\inserted{(13)},
such that $\underline{C}_\infty$ should be the limit (i.e. every arrow in $C_\infty$ should belong \deleted{\illegible{1}}
to some $\underline{C}_n$),  and such that the passage from $\underline{C}_n$ to $\underline{C}_{n+1}$ should satisfy \inserted{(B).}
\deleted{(B).} \hand{$\ulcorner$}It is fairly obvious that these conditions alone do by no means
caracterise $\underline{C}_\infty$  up to equivalence,  and still less the sequence of its
subcategories $\underline{C}_n$ .  The point I wish to make though,  before pursuing
with a proposal of an explicit description,  is \inserted{(firstly)} that \uline{this ambiguity is}
\uline{in the nature of things}\over{,}{.} \over{r}{R}oughly saying, two different mathematicians,
working independently on the conceptual problem I had in mind,  \deleted{\unsure{assuming}}\inserted{assuming}
they both wind up with some explicit definition,  will almost certainly
get non e\over{\illegible{1}}{q}uivalent definitions  -  namely with non equivalent categories
of (set-valued, say) $\infty$-groupoids !  And, secondly and as importantly,

\mspage{29}
